\documentclass{plantillaPPT}

\titulo{11}{Área Encerradas por Curvas Polares}

\begin{document}
\primeraDiapo

\begin{frame}{Descripción e Integración con Curvas Polares}
\framesubtitle{Práctica 11}

1. Sean \(C_1\) y \(C_2\) dos curvas en coordenadas polares definidas por:
\[C_1: r=3\sin(\theta), \; \theta\in[0, \pi] \quad \text{y} \quad C_2:r=1+\sin(\theta), \; \theta\in[0, 2\pi]\]
\((a)\) Escribir en coordenadas rectangulares la región que representa a \(C_1\).\\[2ex]
\((b)\) Verificar que los puntos de intersección entre ambas curvas son \(P_1(\frac{3}{2}, \frac{\pi}{6})\) y \(P_2(\frac{3}{2}, \frac{5\pi}{6})\).\[2ex\]
    
\end{frame}

\begin{frame}{Descripción e Integración con Curvas Polares}
\framesubtitle{Práctica 11}

\((c)\) Escribir la o las integrales que permiten calcular el área de la región achurada para cada caso:

\begin{center}
    \includegraphics[width=1\textwidth]{C2 2025-1/imagenes/ayud11img.png}
\end{center}

\end{frame}

\begin{frame}{Área encerrada por curvas polares}
\framesubtitle{Práctica 11}

2. La figura muestra la gráfica polar \(r = -\cos(\frac{3}{2}\theta)\) para \(0\leq\theta\leq2\pi\). Dar una expresión para determinar el área indicada en la figura, sin desarrollar las integrales.

\begin{center}
    \includegraphics[width=0.5\textwidth]{C2 2025-1/imagenes/ayud11img2.png}
\end{center}
    
\end{frame}

\begin{frame}{Longitud de Arco en Coordenadas Polares}
\framesubtitle{Definición}

\begin{definicion}
    Sea \(f\) una función para la cual \(f'\) es continua sobre un intervalo \(]\alpha, \beta[\). Entonces la longitud \(L\) de la gráfica \(r=f(\theta)\) sobre el intervalo está dada por:

    \[L(C) = \int_\alpha^\beta \sqrt{r^2+\left(\frac{dr}{d\theta}\right)^2}\, d\theta\]
    
\end{definicion}

\end{frame}

\begin{frame}{Cálculo de Longitud de Arco}
\framesubtitle{Práctica 11}

3. Determine la longitud de las siguientes curvas. \\[5ex]

\((a)\) El cardioide \(r=1+\cos(\theta)\), con \(0\leq\theta\leq2\pi\).\\[3ex]
\((b)\) El espiral \(r = e^\theta\), con \(\theta\leq0\). \textbf{Indicación:} Utilice la definición de integral impropia.

\end{frame}

\end{document}